\medskip

%On donne un programme de calcul :
%
%\begin{center}
%\begin{pspicture}(-3,0)(3,4.2)
%%\psgrid
%\rput(0,4){\fbox{choisir un nombre}}\psline[linewidth=1.5pt]{->}(0,3.7)(0,3.2)
%\rput(0,3){\fbox{le multiplier par 2}}\psline[linewidth=1.5pt]{->}(0,2.7)(0,2.2)
%\rput(0,2){\fbox{élever le résultat au carré}}\psline[linewidth=1.5pt]{->}(0,1.7)(0,1.2)
%\rput(0,1){\fbox{retrancher 9}}\psline[linewidth=1.5pt]{->}(0,0.7)(0,0.2)
%\rput(0,0){\fbox{afficher le résultat}}
%\end{pspicture}
%\end{center}

\begin{enumerate}
\item %Lorsque le nombre choisi est 4, vérifier le programme affiche 55, en précisant chacune des
%étapes de calcul.
On a successivement : $4 \longmapsto 8 \longmapsto 64 \longmapsto 55$.
\item On appelle $x$ le nombre choisi au départ.
	\begin{enumerate}
		\item %Écrire, en fonction de $x$, le résultat obtenu par le programme.
		En partant de $x$, on obtient successivement :
		
		$x \longmapsto 2x \longmapsto (2x)^2 = 2^2 \times x^2 = 4x^2 \longmapsto 4x^2 - 9$.
		\item %Parmi les quatre expressions suivantes, laquelle correspond au résultat obtenu par le
%programme ?
	\end{enumerate}
%\begin{center}
%\begin{tabularx}{\linewidth}{*{4}{X}}
%$A = 55$ &$B= (2x +3)^2$ &$C= (2x - 3)(2x +3)$ &$D = (2x — 3)^2$\\
%\end{tabularx}
%\end{center}
On reconnait en $4x^2 - 9$ une différence de deux carrés qui peut se factoriser :

$4x^2 - 9 = (2x)^2 - 3^2 = (2x + 3)(2x - 3) = (2x - 3)(2x + 3)$ : c'est l'expression $C$.
\end{enumerate}

\bigskip
